This project investigates the construction of task-aware knowledge graphs (KGs) from text,
treating graph structure as an optimization target rather than a fixed representation.
Existing frameworks such as Neo4j-based pipelines and LLM-driven tools (e.g., LlamaIndex)
focus on constructing large, general-purpose graphs without incorporating feedback from
downstream tasks, despite evidence from evaluation frameworks (e.g., KGrEaT, GEval, KGBench)
that KG utility is task-dependent. The research addresses this gap by proposing a
feedback-driven approach in which ontology is treated as a tunable component influencing
predictive performance. The study uses document-based datasets linking text to measurable
outcomes, including financial news (FNSPID) paired with stock price movements. LLM agents
are employed for ontology evolution and triple extraction, enabling scalable generation of
multiple KG variants. Graph-derived features are then computed over temporal windows and
used in predictive models, whose performance is fed back to the agents to refine ontology
design, closing the loop between graph construction and downstream task. Preliminary results
indicate that graph structure strongly influences predictive outcomes, with adaptive ontology
updates leading to significant improvements.

% Short version (for submission systems with word limits):
% This project explores building task-aware knowledge graphs (KGs) from text by treating
% graph structure as an optimization target rather than a fixed design. Unlike existing
% approaches that generate large, general-purpose graphs, it introduces a feedback-driven
% framework where ontology is iteratively refined based on downstream task performance.
% Using document-based datasets (e.g., financial news linked to stock movements), LLM agents
% generate and evolve ontologies and extract triples to produce multiple KG variants. Features
% derived from these graphs are used in predictive models, and their performance is fed back
% to guide further ontology updates. Preliminary results show that graph structure
% significantly impacts predictive performance, with adaptive ontology refinement leading to
% measurable improvements.
